\section{Diagonal damping at a fixed Jacobian}
\label{sec:fixedJ-independent}

For a real matrix $J\in\mathbb R^{n\times n}$ and a set
$I\subseteq[n]$, define
\[
  c_I(J)=(-1)^{|I|}\det J[I,I],\qquad c_\varnothing(J)=1.
\]
No stability assumption is made in the next theorem.

\begin{theorem}[Positive diagonal damping]
\label{thm:fixedJ-independent}
The following statements are equivalent.
\begin{enumerate}
\item There is a diagonal matrix $D\succ0$ for which $J-D$ has a
strictly positive real eigenvalue.
\item $c_I(J)<0$ for at least one $I\subseteq[n]$.
\end{enumerate}
When $J$ is Hurwitz, every set in item 2 is nonempty and proper.
\end{theorem}

\begin{proof}
Write $D(d)=\operatorname{diag}(d_1,\ldots,d_n)$.  Multilinearity of the
determinant in its columns gives
\begin{equation}
 \det(D(d)-J)
 =\sum_{I\subseteq[n]}(-1)^{|I|}\det J[I,I]
       \prod_{j\notin I}d_j.
 \label{eq:principal-expansion-independent}
\end{equation}
Indeed, choosing the diagonal column from $D(d)$ at every index in
$I^c$ leaves the principal block $-J[I,I]$ after the corresponding rows
and columns are deleted.  The surviving sign is $(-1)^{|I|}$.  This also
covers $I=\varnothing$, whose contribution is $\prod_jd_j$, and
$I=[n]$, whose contribution is $\det(-J)$.

Suppose first that every $c_I(J)$ is nonnegative.  For arbitrary $D\succ0$
and $\lambda>0$, apply \eqref{eq:principal-expansion-independent} to
$d_i+\lambda$.  Every summand is nonnegative and the empty-set summand is
strictly positive.  Therefore
\[
 \det\bigl(\lambda I-(J-D)\bigr)
 =\det\bigl(D+\lambda I-J\bigr)>0.
\]
A positive real eigenvalue would make this determinant zero, so item 1 is
impossible.  This proves the necessity of item 2 by contraposition.

Now assume $c_I(J)<0$.  For $t>1$ set
\[
 d_j(t)=\begin{cases}
 t,&j\notin I,\\
 t^{-1},&j\in I.
 \end{cases}
\]
The power of $t$ in the monomial indexed by $K$ in
\eqref{eq:principal-expansion-independent} is
\[
 e_I(K)=|(I^c\cap K^c)|-|I\cap K^c|.
\]
For the chosen set $I$, $e_I(I)=n-|I|$, and direct cancellation gives
\[
 e_I(I)-e_I(K)=|I\triangle K|.
\]
Thus the $I$-term is the unique term of highest weight.  Dividing by
$t^{n-|I|}$ and sending $t$ to infinity yields
\[
 t^{-(n-|I|)}\det(D(d(t))-J)\longrightarrow c_I(J)<0.
\]
Hence some finite $t$ gives $\det(D(d(t))-J)<0$, with every diagonal
entry of $D(d(t))$ strictly positive.  The characteristic polynomial
\[
 p(\lambda)=\det\bigl(\lambda I-(J-D(d(t)))\bigr)
\]
is monic, is negative at $\lambda=0$, and tends to $+\infty$ as
$\lambda\to+\infty$.  The intermediate value theorem therefore supplies
a positive real zero.

If $J$ is Hurwitz, then $c_\varnothing=1$.  Also
$c_{[n]}=\det(-J)>0$, because it is the constant coefficient of the monic
Hurwitz characteristic polynomial.  Neither endpoint set can witness item 2.
The argument did not assume full rank; a zero full determinant or any other
degenerate principal minor merely contributes a zero coefficient.
\end{proof}

\begin{remark}[Constructive witness]
The proof is effective over $\mathbb Q$.  Once a negative signed principal
minor is specified, test $t=2,4,8,\ldots$ until the exact determinant becomes
negative.  Termination follows from the unique-leading-term limit.
\end{remark}

\subsection{Designated mobile coordinates}
Let $M\subseteq[n]$ and let $D_M(d)$ have positive diagonal entries on $M$
and zeros on $M^c$.

\begin{theorem}[Positive-spectral-parameter mobile criterion]
\label{thm:mobile-independent}
Assume
\[
 \det(\lambda I-J)>0\quad\text{for every }\lambda>0.
 \label{eq:mobile-background-positive}
\]
Then a matrix $D_M(d)$ exists for which $J-D_M(d)$ has a positive real
eigenvalue if and only if there are $\lambda>0$ and
$I\supseteq M^c$ such that
\[
 \det(\lambda I_{|I|}-J[I,I])<0.
 \label{eq:mobile-block-independent}
\]
\end{theorem}

\begin{proof}
For fixed $\lambda>0$, expansion only in the mobile diagonal variables gives
\begin{equation}
 \det(\lambda I-J+D_M(d))
 =\sum_{I\supseteq M^c}
 \det(\lambda I_{|I|}-J[I,I])
 \prod_{j\in M\setminus I}d_j.
 \label{eq:mobile-expansion-independent}
\end{equation}
If the left side vanishes, the coefficient indexed by $[n]$ is positive by
\eqref{eq:mobile-background-positive}; therefore at least one other
coefficient must be negative.

Conversely, fix a negative coefficient indexed by $I$.  Assign $d_j=t$ for
$j\in M\setminus I$ and $d_j=t^{-1}$ for $j\in M\cap I$.  The exponent gap
between the $I$-monomial and the $K$-monomial is
\[
 |(M\setminus I)\cap K|+|(M\cap I)\setminus K|,
\]
which vanishes only for $K=I$ among sets containing $M^c$.  Hence the chosen
coefficient uniquely dominates and makes
\eqref{eq:mobile-expansion-independent} negative for a finite $t$.  Viewed as
the characteristic polynomial at a new spectral variable $\mu$, the same
determinant is negative at $\mu=\lambda$ and positive for sufficiently large
$\mu$.  It therefore has a real zero $\mu>\lambda>0$.
\end{proof}

The substitution $\lambda=0$ is not valid in general.  For
\[
 J=\begin{pmatrix}-10&20&-50\\1&1&0\\1&0&2\end{pmatrix},
 \qquad M=\{1\},
\]
the matrix $J$ is Hurwitz, the signed determinant of the only proper block
containing the immobile set $\{2,3\}$ is positive, but
\[
 J-\operatorname{diag}(257/2,0,0)
\]
has eigenvalue $3/2$.  At that positive spectral value,
\[
 \det\!\left(\tfrac32I_2-J[\{2,3\},\{2,3\}]\right)=-\tfrac14.
\]
This example separates Theorem~\ref{thm:mobile-independent} from any
zero-parameter principal-minor rule.
